% TOP_PROBLEM: {"id": 2568, "title": "Kourovka Notebook Problem 21.59", "category_id": 17, "category": {"id": 17, "name": "group_theory", "display_name": "Group Theory", "description": "Problems about groups, group actions, representations, and related algebraic structures.", "slug": "group-theory", "order_index": 17, "created_at": "2026-07-31T15:26:25.671Z"}, "status": "open", "problem_number": "KOU-21.59", "set_id": 10, "statusQualification": "Part (b), the quasisimple case, is already established in the cited literature; both original parts are defined, with part (a) retained as the general question."}
% FIRST_POSTED: 2026-10-01
% ENTRY_KIND: statement_only
% UnsolvedMath ID 2568; source catalogue code KOU-21.59
% !TeX program = lualatex
% Definition and sources only; this file is not an LLM solution attempt.
\documentclass[11pt,a4paper]{article}
\usepackage[margin=25mm]{geometry}
\usepackage{fontspec}
\setmainfont{lmroman10-regular.otf}[BoldFont=lmroman10-bold.otf,
 ItalicFont=lmroman10-italic.otf,BoldItalicFont=lmroman10-bolditalic.otf]
\usepackage{amsmath,amssymb,mathtools}
\usepackage{unicode-math}
\setmathfont{latinmodern-math.otf}
\AtBeginDocument{\let\setminus\smallsetminus}
\usepackage{xurl,hyperref}
\hypersetup{colorlinks=true,urlcolor=blue}
\setcounter{secnumdepth}{0}
\setlength{\parindent}{0pt}
\setlength{\parskip}{5pt}
\setlength{\emergencystretch}{3em}
\begin{document}
\section*{Kourovka Notebook Problem 21.59}\label{2568}
{\small UnsolvedMath ID 2568 · KOU-21.59 · Group Theory · Kourovka Notebook - New Problems, Issue 21}
\subsection{Definitions and mathematical statement}
For a finite group $G$, a complex representation is a homomorphism
$G\to\operatorname{GL}(V)$ on a finite-dimensional complex vector space.
It is irreducible if it has no nonzero proper invariant subspace.
Let $\operatorname{Irr}(G)$ index the isomorphism classes of irreducible
representations, and let $\chi(1)=\dim V$ denote the degree of the
corresponding character. Define the multiset
\[
 \chi_1(G)=\{\!\{\chi(1):\chi\in\operatorname{Irr}(G)\}\!\}.
\]
Different irreducibles of the same dimension contribute repeated entries.

Assume finite groups $G,H$ satisfy $\chi_1(G)=\chi_1(H)$.
Part (a) asks whether $H\cong G$ whenever $G$ is almost simple: there
is a nonabelian finite simple group $S$ with
$S\le G\le\operatorname{Aut}(S)$, identifying $S$ with its inner
automorphisms. Part (b) asks the same when $G$ is quasisimple, meaning
$G=G'$ and $G/Z(G)$ is nonabelian simple. Here $G'$ is the commutator
subgroup and $Z(G)$ the subgroup commuting with all elements.

The comparison group $H$ is arbitrary; it is not assumed almost simple
or quasisimple. Equality concerns the degree multiset, not full character
tables or only the set of distinct degrees. In particular it implies
equality of group orders through the sum of the squares of the degrees.

Both catalogue parts are retained. Part (b) is already recorded as a
theorem in the cited primary literature; the remaining general recognition
question in this entry is part (a). No new proof of that settled part
is claimed here.
\subsection{Short English statement}
Do irreducible complex character degrees, counted with multiplicity, determine an almost simple group or a quasisimple group up to isomorphism?
\subsection{Sources}
\begin{itemize}
\item[S1] ULAM.AI and the UnsolvedMath Contributors, supplied problems.json, record 2568 (KOU-21.59), Kourovka Notebook - New Problems, Issue 21. Catalogue identity and source transcription; CC BY 4.0.
\url{https://huggingface.co/datasets/ulamai/UnsolvedMath}.
\item[S2] The Kourovka Notebook, 21st edition, new problems, Problem 21.59. Expanded formulation of the supplied catalogue statement.
\url{https://arxiv.org/pdf/1401.0300v47}.
\item[S3] A. Moretó, Multiplicities of character codegrees of finite groups, Bulletin of the London Mathematical Society (2023), introduction records the quasisimple recognition theorem.
\url{https://doi.org/10.1112/blms.12724}.
\end{itemize}
\end{document}
